\documentclass[10pt]{extarticle}
\usepackage[a4paper,margin=1.5cm]{geometry}
\usepackage[T1]{fontenc}
\usepackage{amsmath,amssymb}
\usepackage{graphicx}
\usepackage[export]{adjustbox}
\usepackage{multicol}
\usepackage{enumitem}
\usepackage{booktabs,array,tabularx}
\setlength{\columnsep}{0.7cm}
\setlength{\parindent}{0pt}
\setlength{\parskip}{3pt}
\pagestyle{empty}
\makeatletter
% figure registry (figdefs.tex is generated by the extract step)
\newcommand{\qfigdef}[4]{\@namedef{qf@file@#1}{#2}\@namedef{qf@w@#1}{#3}\@namedef{qf@wide@#1}{#4}}
\newcommand{\qf@check}[1]{\@ifundefined{qf@file@#1}{\PackageError{score2}{Unknown figure id #1}{}}{}}
\newcommand{\qf@img}[4][]{\includegraphics[width=#2\dimexpr\@nameuse{qf@w@#4}pt\relax,max width=#3,max height=0.45\textheight,#1]{\@nameuse{qf@file@#4}}}
% Full-width questions: finalize wraps every question that contains a wide figure in \qfullbegin...\qfullend,
% so the WHOLE question (stem, figures, options) is set across the page in reading order, then the two columns
% resume. The question's head - its number, stem and first figure (\qfullsplit marks the end of the head) -
% is kept on one page, so a number is never left at a page foot with its figure overleaf; the rest may break
% across pages normally, so no big gaps are left. Without a \qfullsplit, the whole question is kept together
% if it is under 45% of a page. Never write these yourself.
\newif\ifqf@full
\newif\ifqf@boxed
\newsavebox{\qf@box}
\newcommand{\qf@open}{\global\qf@fulltrue\par\global\qf@boxedtrue
  \setbox\qf@box\vbox\bgroup\hsize\textwidth\linewidth\textwidth}
\newcommand{\qf@emit}[1]{\par\egroup\global\qf@boxedfalse
  \ifdim\dimexpr\ht\qf@box+\dp\qf@box\relax<#1\textheight\noindent\usebox{\qf@box}\par\else\unvbox\qf@box\fi}
\newcommand{\qfullsplit}{\ifqf@boxed\qf@emit{0.85}\fi}
\newcommand{\qf@close}{\ifqf@boxed\qf@emit{0.45}\fi\global\qf@fullfalse}
\newcommand{\qfullbegin}{\end{multicols}\qf@open}
\newcommand{\qfullend}{\qf@close\begin{multicols}{2}}
\newcommand{\qfullstart}{\qf@open}
\newcommand{\qfulldone}{\qf@close}
% \qfig[scale]{id}: block figure, centred, directly where it is written (size set by the tool).
% Fallback if a wide figure is not inside a full-width question: break out of the columns in place - no floats.
\newcommand{\qfig}[2][1]{\qf@check{#2}%
  \ifqf@full\qf@blk{#1}{#2}\else
  \ifnum\@nameuse{qf@wide@#2}=1 \expandafter\@firstoftwo\else\expandafter\@secondoftwo\fi
  {\end{multicols}\begin{center}\qf@img{#1}{\textwidth}{#2}\end{center}\begin{multicols}{2}}%
  {\qf@blk{#1}{#2}}\fi}
\newcommand{\qf@blk}[2]{\par\smallskip{\centering\qf@img{#1}{\linewidth}{#2}\par}\smallskip}
% \qfigi[scale]{id}: inline figure (inside option lists / tables / minipages); never breaks columns.
\newcommand{\qfigi}[2][1]{\qf@check{#2}\qf@img[valign=c]{#1}{\linewidth}{#2}}
\makeatother
% structure macros - lint reads these to check completeness and figure placement
\newcommand{\papersec}[1]{\par\bigskip{\centering\Large\bfseries PAPER-#1\par}\medskip}
\newcommand{\partsec}[2]{\par\medskip\noindent{\large\bfseries PART-#1: #2}\par\smallskip}
\newcommand{\sect}[1]{\par\smallskip\noindent\textbf{#1}\par}
\newcommand{\qn}[1]{\par\smallskip\noindent\textbf{#1.}~}
\newcommand{\paraq}[2]{\par\smallskip\noindent\textbf{Paragraph for Questions #1 and #2}\par}
\newenvironment{qoptions}{\begin{itemize}[leftmargin=1.8em,itemsep=1pt,topsep=2pt,label={}]}{\end{itemize}}
\newcommand{\opt}[2]{\item[(#1)] #2}
% figure registry (generated by score2_tools.py finalize)
\qfigdef{p1_phy_s1_q01_1}{Score Advanced - Test 08 (22 Apr 2023)-img/p1_phy_s1_q01_1.png}{148.0}{0}
\qfigdef{p1_phy_s1_q02_1}{Score Advanced - Test 08 (22 Apr 2023)-img/p1_phy_s1_q02_1.png}{149.4}{0}
\qfigdef{p1_phy_s1_q04_1}{Score Advanced - Test 08 (22 Apr 2023)-img/p1_phy_s1_q04_1.png}{128.9}{0}
\qfigdef{p1_phy_s1_q06_1}{Score Advanced - Test 08 (22 Apr 2023)-img/p1_phy_s1_q06_1.png}{61.3}{0}
\qfigdef{p1_phy_s2_q08_1}{Score Advanced - Test 08 (22 Apr 2023)-img/p1_phy_s2_q08_1.png}{195.9}{0}
\qfigdef{p1_phy_s2_q09_1}{Score Advanced - Test 08 (22 Apr 2023)-img/p1_phy_s2_q09_1.png}{190.0}{0}
\qfigdef{p1_phy_s2_q11_1}{Score Advanced - Test 08 (22 Apr 2023)-img/p1_phy_s2_q11_1.png}{92.5}{0}
\qfigdef{p1_phy_s3_q02_1}{Score Advanced - Test 08 (22 Apr 2023)-img/p1_phy_s3_q02_1.png}{140.1}{0}
\qfigdef{p1_phy_s3_q03_1}{Score Advanced - Test 08 (22 Apr 2023)-img/p1_phy_s3_q03_1.png}{122.4}{0}
\qfigdef{p1_phy_s3_q05_1}{Score Advanced - Test 08 (22 Apr 2023)-img/p1_phy_s3_q05_1.png}{137.9}{0}
\qfigdef{p1_phy_s3_q06_2}{Score Advanced - Test 08 (22 Apr 2023)-img/p1_phy_s3_q06_2.png}{220.6}{0}
\qfigdef{p1_chem_s1_q02_A}{Score Advanced - Test 08 (22 Apr 2023)-img/p1_chem_s1_q02_A.png}{133.7}{0}
\qfigdef{p1_chem_s1_q02_B}{Score Advanced - Test 08 (22 Apr 2023)-img/p1_chem_s1_q02_B.png}{133.0}{0}
\qfigdef{p1_chem_s1_q02_C}{Score Advanced - Test 08 (22 Apr 2023)-img/p1_chem_s1_q02_C.png}{127.5}{0}
\qfigdef{p1_chem_s1_q02_D}{Score Advanced - Test 08 (22 Apr 2023)-img/p1_chem_s1_q02_D.png}{135.4}{0}
\qfigdef{p1_chem_s2_q09_B}{Score Advanced - Test 08 (22 Apr 2023)-img/p1_chem_s2_q09_B.png}{155.8}{0}
\qfigdef{p1_chem_s2_q09_C}{Score Advanced - Test 08 (22 Apr 2023)-img/p1_chem_s2_q09_C.png}{184.2}{0}
\qfigdef{p1_chem_s2_q09_D}{Score Advanced - Test 08 (22 Apr 2023)-img/p1_chem_s2_q09_D.png}{148.2}{0}
\qfigdef{p1_chem_s2_q10_1}{Score Advanced - Test 08 (22 Apr 2023)-img/p1_chem_s2_q10_1.png}{125.8}{0}
\qfigdef{p1_chem_s2_q12_A}{Score Advanced - Test 08 (22 Apr 2023)-img/p1_chem_s2_q12_A.png}{215.1}{0}
\qfigdef{p1_chem_s2_q12_D}{Score Advanced - Test 08 (22 Apr 2023)-img/p1_chem_s2_q12_D.png}{245.2}{0}
\qfigdef{p1_chem_s3_q03_1}{Score Advanced - Test 08 (22 Apr 2023)-img/p1_chem_s3_q03_1.png}{110.2}{0}
\qfigdef{p1_chem_s3_q04_1}{Score Advanced - Test 08 (22 Apr 2023)-img/p1_chem_s3_q04_1.png}{102.4}{0}
\qfigdef{p1_chem_s3_q06_2}{Score Advanced - Test 08 (22 Apr 2023)-img/p1_chem_s3_q06_2.png}{40.4}{0}
\qfigdef{p1_chem_s3_q06_4}{Score Advanced - Test 08 (22 Apr 2023)-img/p1_chem_s3_q06_4.png}{40.2}{0}
\qfigdef{p1_chem_s3_q06_3}{Score Advanced - Test 08 (22 Apr 2023)-img/p1_chem_s3_q06_3.png}{57.3}{0}
\qfigdef{p1_chem_s3_q06_1}{Score Advanced - Test 08 (22 Apr 2023)-img/p1_chem_s3_q06_1.png}{38.5}{0}
\qfigdef{p1_chem_s3_q06_7}{Score Advanced - Test 08 (22 Apr 2023)-img/p1_chem_s3_q06_7.png}{40.6}{0}
\qfigdef{p1_chem_s3_q06_5}{Score Advanced - Test 08 (22 Apr 2023)-img/p1_chem_s3_q06_5.png}{55.9}{0}
\qfigdef{p1_chem_s3_q06_6}{Score Advanced - Test 08 (22 Apr 2023)-img/p1_chem_s3_q06_6.png}{38.4}{0}
\qfigdef{p1_chem_s3_q06_8}{Score Advanced - Test 08 (22 Apr 2023)-img/p1_chem_s3_q06_8.png}{46.8}{0}
\qfigdef{p2_phy_s1_q01_1}{Score Advanced - Test 08 (22 Apr 2023)-img/p2_phy_s1_q01_1.png}{73.7}{0}
\qfigdef{p2_phy_s1_q02_1}{Score Advanced - Test 08 (22 Apr 2023)-img/p2_phy_s1_q02_1.png}{145.8}{0}
\qfigdef{p2_phy_s1_q03_1}{Score Advanced - Test 08 (22 Apr 2023)-img/p2_phy_s1_q03_1.png}{152.6}{0}
\qfigdef{p2_phy_s1_q04_1}{Score Advanced - Test 08 (22 Apr 2023)-img/p2_phy_s1_q04_1.png}{139.8}{0}
\qfigdef{p2_phy_s1_q05_1}{Score Advanced - Test 08 (22 Apr 2023)-img/p2_phy_s1_q05_1.png}{207.2}{0}
\qfigdef{p2_phy_s1_q06_1}{Score Advanced - Test 08 (22 Apr 2023)-img/p2_phy_s1_q06_1.png}{105.5}{0}
\qfigdef{p2_phy_s2_q01_1}{Score Advanced - Test 08 (22 Apr 2023)-img/p2_phy_s2_q01_1.png}{162.7}{0}
\qfigdef{p2_phy_s2_q03_1}{Score Advanced - Test 08 (22 Apr 2023)-img/p2_phy_s2_q03_1.png}{131.7}{0}
\qfigdef{p2_phy_s2_q04_1}{Score Advanced - Test 08 (22 Apr 2023)-img/p2_phy_s2_q04_1.png}{122.5}{0}
\qfigdef{p2_phy_s3_q01_1}{Score Advanced - Test 08 (22 Apr 2023)-img/p2_phy_s3_q01_1.png}{147.0}{0}
\qfigdef{p2_phy_s3_q03_1}{Score Advanced - Test 08 (22 Apr 2023)-img/p2_phy_s3_q03_1.png}{117.6}{0}
\qfigdef{p2_chem_s1_q01_1}{Score Advanced - Test 08 (22 Apr 2023)-img/p2_chem_s1_q01_1.png}{133.4}{0}
\qfigdef{p2_chem_s1_q02_1}{Score Advanced - Test 08 (22 Apr 2023)-img/p2_chem_s1_q02_1.png}{172.5}{0}
\qfigdef{p2_chem_s1_q05_1}{Score Advanced - Test 08 (22 Apr 2023)-img/p2_chem_s1_q05_1.png}{60.5}{0}
\qfigdef{p2_chem_s1_q05_A}{Score Advanced - Test 08 (22 Apr 2023)-img/p2_chem_s1_q05_A.png}{24.2}{0}
\qfigdef{p2_chem_s1_q05_B}{Score Advanced - Test 08 (22 Apr 2023)-img/p2_chem_s1_q05_B.png}{67.1}{0}
\qfigdef{p2_chem_s1_q05_C}{Score Advanced - Test 08 (22 Apr 2023)-img/p2_chem_s1_q05_C.png}{28.9}{0}
\qfigdef{p2_chem_s1_q05_D}{Score Advanced - Test 08 (22 Apr 2023)-img/p2_chem_s1_q05_D.png}{32.5}{0}
\qfigdef{p2_chem_s2_q06_1}{Score Advanced - Test 08 (22 Apr 2023)-img/p2_chem_s2_q06_1.png}{156.5}{0}
\qfigdef{p2_chem_s3_q03_2}{Score Advanced - Test 08 (22 Apr 2023)-img/p2_chem_s3_q03_2.png}{157.5}{0}
\qfigdef{p2_chem_s3_q05_2}{Score Advanced - Test 08 (22 Apr 2023)-img/p2_chem_s3_q05_2.png}{69.8}{0}
\qfigdef{p2_chem_s3_q05_1}{Score Advanced - Test 08 (22 Apr 2023)-img/p2_chem_s3_q05_1.png}{120.8}{0}
\qfigdef{p2_chem_s3_q05_3}{Score Advanced - Test 08 (22 Apr 2023)-img/p2_chem_s3_q05_3.png}{198.7}{0}
\qfigdef{p2_math_s3_q02_1}{Score Advanced - Test 08 (22 Apr 2023)-img/p2_math_s3_q02_1.png}{70.5}{0}
\raggedcolumns
\makeatletter\@beginparpenalty=10000\makeatother
\newcolumntype{L}[1]{>{\raggedright\arraybackslash}p{#1}}
\begin{document}
\begin{center}{\Large\bfseries Score Advanced --- Test 08 (22 Apr 2023)}\end{center}
\begin{multicols}{2}
\raggedcolumns
\papersec{1}
\partsec{1}{PHYSICS}
\sect{SECTION-I (i)}
\qn{1} Three small discs are connected with two identical massless rods as shown in figure. The rods are pinned to the discs such that angle between them can change freely. The system is placed on a smooth horizontal surface with discs P and Q touching a smooth wall. A force F is applied to the disc R is in direction perpendicular to the wall. Find acceleration of disc P immediately after the force starts to act. Masses of discs are m$_\text{P}$ = m$_\text{Q}$ = 2 m, m$_\text{R}$ = m [wall is perpendicular to the plane of the figure.]\par
\qfig{p1_phy_s1_q01_1}
\begin{qoptions}
\opt{A}{$\dfrac{\text{F}}{\sqrt{2}\text{m}}$}
\opt{B}{$\dfrac{\text{F}}{8\text{m}}$}
\opt{C}{$\dfrac{2\text{F}}{11\text{m}}$}
\opt{D}{None of these}
\end{qoptions}
\qn{2} The circuit shown, consists of three resistances, an incandescent bulb, a switch and an ideal battery. If the bulb B glows equally bright whether the switch is closed or open and if the resistance of the bulb is 10 P $\Omega$, then find the value of P.\par
\qfig{p1_phy_s1_q02_1}
\begin{qoptions}
\opt{A}{0}
\opt{B}{1}
\opt{C}{3}
\opt{D}{5}
\end{qoptions}
\qn{3} Two identical particles A and B of mass m each are connected together by a light and inextensible string of length $\ell$. The particles are held at rest in air in same horizontal level at a separation $\ell$. Both particles are released simultaneously and one of them (say A) is given speed v$_0$ vertically upward. Choose the correct option. Ignore air resistance and gravity to be uniform and particles not collide with ground.
\begin{qoptions}
\opt{A}{The maximum height attained by the centre of mass of the system of A and B is dependent on $\ell$}
\opt{B}{The kinetic energy of the system of A and B when the centre of mass is at its highest point, is $\dfrac{\text{mv}_0^2}{8}$}
\opt{C}{The maximum height attained by the centre of mass of the system of A and B is $\dfrac{\text{v}_0^2}{4\text{g}}$}
\opt{D}{The kinetic energy of the system of A and B when the centre of mass is at its highest point, is $\dfrac{\text{mv}_0^2}{4}$}
\end{qoptions}
\qn{4} In a gravity free space, three identical metal plates have been placed parallel to each other as shown. In the initial state shown in the diagram, the potential difference between the points A and B is V$_1$. Now, the middle plate is shifted down a distance L/2, keeping the battery connected. The other plates are kept fixed. The potential difference between the points A and B now becomes V$_2$. The value V$_1$/V$_2$ is equal to\par
\qfig{p1_phy_s1_q04_1}
\begin{qoptions}
\opt{A}{1}
\opt{B}{5/6}
\opt{C}{5/4}
\opt{D}{3/2}
\end{qoptions}
\qn{5} A particle executing SHM with time period T and amplitude A along horizontal direction along x axis. When it is at +$\dfrac{\text{A}}{2}$ from equilibrium position and moving towards right, it received an impulse which doubles its velocity without changing its direction. The time taken by particle to reach extreme right position is $\dfrac{\text{T}}{2\pi}\cos^{-1}\left[\dfrac{1}{\text{K}}\right]$. The value of K is :-
\begin{qoptions}
\opt{A}{$\sqrt{8}$}
\opt{B}{$\sqrt{12}$}
\opt{C}{$\sqrt{13}$}
\opt{D}{$\sqrt{15}$}
\end{qoptions}
\qn{6} A wheel of radius R is rolling without sliding inside a fixed circular cylinder of radius 2R as shown. What is the trajectory followed by a point on the rim of the wheel ?\par
\qfig{p1_phy_s1_q06_1}
\begin{qoptions}
\opt{A}{Circle}
\opt{B}{Cycloid}
\opt{C}{Straight Line}
\opt{D}{None of these}
\end{qoptions}
\sect{SECTION-I (ii)}
\qn{7} A car of mass m is moving on a horizontal surface with speed v$_0$ along positive x-axis. The driver of car wishes to accelerate the car. He accelerates the car which provides a constant tractive force F, directed along the line of motion of car. The frontal force of resistance is directed opposite to the velocity and given as F$_\text{r}$ = $\alpha$v$^2$ in magnitudes where $\alpha$ is a positive constant and v is speed of car at that time. Assume initially car is at origin. Express the tractive force F as a function of x and v, where x is the displacement of the car and v is the velocity at that position.
\begin{qoptions}
\opt{A}{$\alpha\,\dfrac{\text{v}^2-\text{v}_0^2\text{e}^{\frac{-2\alpha\text{x}}{\text{m}}}}{1-\text{e}^{\frac{-2\alpha\text{x}}{\text{m}}}}$}
\opt{B}{$\alpha\,\dfrac{\text{v}^2+\text{v}_0^2\text{e}^{\frac{-2\alpha\text{x}}{\text{m}}}}{1-\text{e}^{\frac{-2\alpha\text{x}}{\text{m}}}}$}
\opt{C}{$\alpha\,\dfrac{\text{v}^2-\text{v}_0^2\text{e}^{\frac{-2\alpha\text{x}}{\text{m}}}}{1+\text{e}^{\frac{-2\alpha\text{x}}{\text{m}}}}$}
\opt{D}{$\alpha\,\dfrac{\text{v}^2+\text{v}_0^2\text{e}^{\frac{-2\alpha\text{x}}{\text{m}}}}{1+\text{e}^{\frac{-2\alpha\text{x}}{\text{m}}}}$}
\end{qoptions}
\qn{8} Consider all pulleys and string to be ideal.\par
\qfig{p1_phy_s2_q08_1}
\begin{qoptions}
\opt{A}{the value of T$_2$ is 200 N}
\opt{B}{the value of T$_2$ is $\dfrac{400}{3}$N}
\opt{C}{the force of friction between B and C is 70 N}
\opt{D}{the force of friction between B and C is 90 N}
\end{qoptions}
\qn{9} In the given circuit the potential difference between a \& b is 300 V. The resistors are of 100 $\Omega$ each.\par
\qfig{p1_phy_s2_q09_1}
\begin{qoptions}
\opt{A}{If the upper meter M$_1$ is ideal voltmeter \& lower is ideal ammeter M$_2$, reading of voltmeter is 300 V.}
\opt{B}{If both meters are ideal voltmeters, they show 200 V each.}
\opt{C}{If both meters are ideal ammeters, they read 6A each.}
\opt{D}{If the upper meter M$_1$ is ideal voltmeter and lower is ideal ammeter M$_2$, reading of ammeter is 3A.}
\end{qoptions}
\qn{10} A steady beam of $\alpha$-particles travelling with kinetic energy E = 83.5 keV carries a current of I = 0.2 $\mu$A. If this beam strikes a plane surface at an angle $\theta$ = 30$^\circ$ with normal to the surface, then choose the correct statement(s) :- (Mass of $\alpha$-particles = 6.68 $\times$ 10$^{-27}$ kg.)
\begin{qoptions}
\opt{A}{The number of $\alpha$-particles strike the surface in t = 4 sec is 2.5 $\times$ 10$^{12}$.}
\opt{B}{The number of $\alpha$-particles are there in length $\ell$ = 20 cm of the beam is 6.25 $\times$ 10$^4$.}
\opt{C}{The power of the source used to accelerate these $\alpha$-particles from rest is 8.35 $\times$ 10$^{-3}$ watts.}
\opt{D}{The power of the source used to accelerate these $\alpha$-particles from rest is 3.23 $\times$ 10$^{-3}$ watts.}
\end{qoptions}
\qn{11} Four identical sphere of radius R/4 are taken out from a uniform sphere of mass M and radius R as shown in the figure. The density of sphere is uniform. Choose the correct option(s).\par
\qfig{p1_phy_s2_q11_1}
\begin{qoptions}
\opt{A}{Gravitational field at the origin is zero.}
\opt{B}{Gravitational potential at the centre of the spherical cavity is $\dfrac{\text{GM}}{64\text{R}}\left(81-2\sqrt{2}\right)$.}
\opt{C}{The work done by gravitational field in moving a point mass from centre of third cavity to centre of second cavity is zero.}
\opt{D}{Gravitational potential at the centre of second cavity is $-\dfrac{\text{GM}}{64\text{R}}\left(75-2\sqrt{2}\right)$.}
\end{qoptions}
\qn{12} A standing wave is set up in a string of variable length and tension by a vibrator of variable frequency. Both ends of the string are fixed. When the vibrator has a frequency f, in a string of length L and under tension T, n antinodes are set up in the string. Mark the correct statement(s) :-
\begin{qoptions}
\opt{A}{If the length of the string is doubled, the frequency should be changed by a factor of 1/2 so that same number of antinodes are produced.}
\opt{B}{If the length of the string is doubled, the frequency should be changed by a factor of 2 so that same number of antinodes are produced.}
\opt{C}{If the frequency and length are held constant, tension needed to produce n + 1 antinodes is given by $\left[\dfrac{\text{n}}{(\text{n}+1)}\right]^2$T.}
\opt{D}{If the frequency is tripled and the length of the string is halved, then tension should be changed by factor of $\dfrac{9}{16}$ so that twice as many antinodes are produced.}
\end{qoptions}
\sect{SECTION-II}
\qn{1} A certain planet of radius R is composed of a uniform material that, through radioactive decay, generates a net power P such that heat per second $\dfrac{\Delta\text{Q}}{\Delta\text{t}}$ = P$\dfrac{\text{r}^3}{\text{R}^3}$, where r is radial distance from centre. This results in a temperature differential between the inside and outside of the planet as heat is transferred from the interior to the surface. The rate of heat transfer inside is governed by conduction. It is found that thermal conductivity k is constant for the planet. For the following assume that the planet is in a steady state; temperature might depend on position, but does not depend on time. (Assuming black body radiation, and emissivity is 1) Find an expression for the temperature difference between the surface of the planet (T$_\text{S}$) and the centre of the planet (T$_\text{C}$). If your answer is $\Delta$T = $\dfrac{\sigma\text{R}^\text{m}\text{T}_\text{S}^\text{n}}{\ell\text{k}}$ fill value of (m + n + $\ell$)
\qn{2} Consider a plane inclined at an angle 45$^\circ$ with the horizontal has two slits (S$_1$ and S$_2$) separated by a distance d = $\sqrt{2}$ mm. The screen is placed at a distance of D = 10 m. A parallel monochromatic light beam of wavelength 5000 \AA\ is incident on the slits as shown. If the fringe width of interference pattern on the screen is k x 10$^{-3}$ meter, then find the value of k.\par
\qfig{p1_phy_s3_q02_1}
\qn{3} A sphere of radius R is kept on smooth horizontal surface. It is connected with ideal spring of constant k as shown in figure. A parallel light beam having intensity I is falling at an angle $\theta$ = 30$^\circ$ with vertical on the sphere, then find the compression in the spring (in SI unit) in equilibrium position, assuming that the sphere absorbs 100\% of the light incident on it. [Take $\dfrac{\text{IR}^2}{\text{KC}}$ = $\dfrac{10}{\pi}$ SI unit, where C is speed of light in vacuum]\par
\qfig{p1_phy_s3_q03_1}
\qn{4} A vernier calliper has no zero error, one main scale division = 1 mm and 9 main scale division = 10 vernier scale divisions. A spring is held between jaws without exerting force on jaws. Main scale reading is 3.4 cm and 5$^{\text{th}}$ vernier scale division coincides with a main division. When Jaws are pressed by force of 10 N on both sides, main scale reading remains 3.4 cm, but 1$^{\text{st}}$ division of vernier scale coincides with a main scale division. Now spring is removed, cut into two pieces and when one part is inserted between jaws without any force, main scale reading = 1.1 cm and 5$^{\text{th}}$ vernier scale division coincides with a main scale. If we apply force of 30 N on both the jaws, main scale reading remains 1.1 cm and x vernier scale division coincides with a main scale division. Find x.
\qn{5} A uniform solid sphere of radius R floats in water (density $\rho_\omega$) such that distance between the water surface and lowest end of the sphere is $\dfrac{3\text{R}}{2}$ as shown in figure. Then the specific gravity of the sphere is......\par
\qfig{p1_phy_s3_q05_1}
\qn{6} A wooden platform can be rotated about its vertical axis with constant small angular speed $\omega$ with the help of a motor. A buzzer is fixed at the circumference of the platform and it rotates in a circle of radius R. The buzzer produces sound of frequency f$_0$. A mic is placed just beneath the platform near its circumference. An electronic frequency analyzer connected to the mic records the frequency (f) received by the mic. Take time (t) to be zero when the buzzer is just above the mic. The frequency f received by mic is found to be f$_0$ $\left[\dfrac{\text{xV}_0}{\text{yV}_0+\text{R}\omega\left[\cos\left(\frac{\omega\text{t}}{\text{z}}\right)\right]}\right]$. The value of (x + y + z) is .........\par
(Speed of sound = V$_0$ and R$\omega$ is not comparable with the speed of sound)\par
\qfig{p1_phy_s3_q06_2}
\partsec{2}{CHEMISTRY}
\sect{SECTION-I (i)}
\qn{1} Which of the following process is correctly related to the extraction of Mg from sea water ?
\begin{qoptions}
\opt{A}{Sea water $\xrightarrow{\text{Na}_2\text{CO}_3}$ MgCO$_3$ $\xrightarrow{\text{Calcination}}$ MgO $\xrightarrow{\text{Carbon}}$ Mg}
\opt{B}{Sea water $\xrightarrow{\text{lime water}}$ Mg(OH)$_2$ $\xrightarrow{\text{HCl(aq.)}}$ MgCl$_2\cdot$6H$_2$O $\xrightarrow[\text{with excess of dry HCl(g)}]{\Delta}$ MgCl$_2$ (Anhyd.) $\xrightarrow{\text{Electrolysis in molten condition}}$ Mg}
\opt{C}{Sea water $\xrightarrow{\text{drying}}$ Sea salt $\xrightarrow{\text{Electrolysis}}$ Mg}
\opt{D}{Sea water $\xrightarrow{\text{Ca(OH)}_2}$ Mg(OH)$_2$ $\xrightarrow{\text{HCl(aq.)}}$ MgCl$_2$ Solution $\xrightarrow{\text{Electrolysis}}$ Mg}
\end{qoptions}
\qn{2} Which of the following is not correct order of reactivity.
\begin{qoptions}
\opt{A}{\qfigi{p1_chem_s1_q02_A}}
\opt{B}{\qfigi{p1_chem_s1_q02_B}}
\opt{C}{\qfigi{p1_chem_s1_q02_C}}
\opt{D}{\qfigi{p1_chem_s1_q02_D}}
\end{qoptions}
\qn{3} Which of the following information is not correct with respect to enzymes and their catalytic action.
\begin{qoptions}
\opt{A}{Enzymes are complex nitrogenous compounds which form colloidal solution in water.}
\opt{B}{Enzymes are protein molecules of high molecular mass where as co-enzymes are non protein used for enhancing catalytic action of enzyme.}
\opt{C}{The rate of enzyme catalysed reaction is higher at high pH of around 12.}
\opt{D}{The enzyme urease cannot catalyze hydrolysis of any amide other than urea.}
\end{qoptions}
\qn{4} Which of the following information is correct with respect to group 15 elements and their compounds.
\begin{qoptions}
\opt{A}{Colourless gas is produced when concentrated HNO$_3$ reacts with Zn and Cu.}
\opt{B}{N--N bond length in N$_2$O$_3$ and N$_2$O$_4$ are almost comparable.}
\opt{C}{Phosphorous acid on heating disproportionates to give phosphoric acid and hypophosphorous acid.}
\opt{D}{White phosphorous on reaction with alkali gives salt of hypophosphorous acid whereas red phosphorous on heating with alkali gives salt of hypophosphoric acid.}
\end{qoptions}
\qn{5} In Carius method of estimation of halogen, 0.15 g of an organic compound gave 0.12 g of AgBr. What is the percentage of bromine in the compound. (Ag = 108, Br = 80)
\begin{qoptions}
\opt{A}{18\%}
\opt{B}{94\%}
\opt{C}{63\%}
\opt{D}{34\%}
\end{qoptions}
\qn{6} For the following reaction : 2A + B + C $\longrightarrow$ 3D, rate was measured at different concentrations and noted. Identify the option which is correct on the basis of data.\par
{\scriptsize\setlength{\tabcolsep}{1pt}\noindent\begin{tabular}{|c|c|c|c|c|}
\hline
Exp. & conc.of [A] & conc. of [B] & conc. of [C] & Rate of formation of [D] \\ \hline
(1) & 0.03 M & 0.04 M & 0.02 M & 10$^{-3}$ M min$^{-1}$ \\ \hline
(2) & 0.06 M & 0.08 M & 0.02 M & 10$^{-3}$ M min$^{-1}$ \\ \hline
(3) & 0.06 M & 0.16 M & 0.04 M & 5 $\times$ 10$^{-4}$ M min$^{-1}$ \\ \hline
(4) & 0.06 M & 0.08 M & 0.04 M & 10$^{-3}$ M min$^{-1}$ \\ \hline
\end{tabular}}\par
\begin{qoptions}
\opt{A}{Overall order of the reaction is negative}
\opt{B}{Order of the reaction with respect to A will be 1.}
\opt{C}{Overall order of the reaction is 2}
\opt{D}{Order of the reaction with respect to B will be 1.}
\end{qoptions}
\sect{SECTION-I (ii)}
\qn{7} An ideal gas with adiabatic component undergoes a process in which internal energy depends upon volume as U = aV$^\alpha$ then select correct statements:
\begin{qoptions}
\opt{A}{Molar heat capacity of process is $\left(\dfrac{\text{R}}{\alpha}-\dfrac{\text{R}}{1-\gamma}\right)$}
\opt{B}{Heat change in process is given by $\Delta$U $\left[1+\dfrac{(\gamma-1)}{\alpha}\right]$}
\opt{C}{Change in internal energy is $\dfrac{(\gamma-1)}{\alpha}\Delta$T}
\opt{D}{Equation of process is PV$^\text{x}$ = constant, where x $>$ 0 if $\alpha$ $<$ 1.}
\end{qoptions}
\qn{8} Which of the following statements is/are incorrect?
\begin{qoptions}
\opt{A}{Correct order of 2$^{\text{nd}}$ I.E. is O $>$ N $>$ C $>$ B}
\opt{B}{Correct order of EA is O $>$ C $>$ B $>$ Li $>$ N}
\opt{C}{Element with Z = 72 belongs to 6$^{\text{th}}$ period \& p-block}
\opt{D}{I.E. order for 13$^{\text{th}}$ group element is B $>$ T$\ell$ $>$ Al $>$ Ga $>$ In}
\end{qoptions}
\qn{9} In which of the following reactions cinnamic acid is one of the major organic product ?
\begin{qoptions}
\opt{A}{Ph--CH$_2$--COOEt $\xrightarrow[\text{(ii) H}_3\text{O}^{\oplus}\text{; (iii) }\Delta]{\text{(i) Et}\bar{\text{O}}\text{Na}^{\oplus}}$}
\opt{B}{\qfigi{p1_chem_s2_q09_B}}
\opt{C}{\qfigi{p1_chem_s2_q09_C}}
\opt{D}{\qfigi{p1_chem_s2_q09_D}}
\end{qoptions}
\qn{10} Which of the following statement(s) is / are correct for the gas taken?\par
\qfig{p1_chem_s2_q10_1}
\begin{qoptions}
\opt{A}{It is H$_2$ or He gas at any temperature.}
\opt{B}{It is any non--ideal gas at temperature above its Boyle's temperature.}
\opt{C}{It is any non--ideal gas at Boyle's temperature.}
\opt{D}{It is H$_2$ or He gas at 0$^\circ$ C.}
\end{qoptions}
\qn{11} Correct statements is/are :
\begin{qoptions}
\opt{A}{SF$_4$ on hydrolysis yield H$_2$SO$_4$.}
\opt{B}{Chlorine gas is yellow green coloured as energy difference between HOMO \& LUMO is corresponding to visible range.}
\opt{C}{In alkali metals, as we move down the group, for there fluorides $\Delta$H$_\text{f}$ values become less negative.}
\opt{D}{LiF is least soluble in alkali metal fluoride due to high lattice energy.}
\end{qoptions}
\qn{12} Which of the following observation is/are correct ?
\begin{qoptions}
\opt{A}{\qfigi{p1_chem_s2_q12_A}}
\opt{B}{C$_6$H$_{14}$ $\xrightarrow[\text{(ii) Oleum; (iii) NaOH/}\Delta]{\text{(i) Al}_2\text{O}_3\text{--Cr}_2\text{O}_3}$ A $\xrightarrow[\text{conc.H}_2\text{SO}_4]{\text{Phthalic Anhydride}}$ $\xrightarrow{\text{NaOH}}$ Pink colour}
\opt{C}{C$_2$H$_6$O (Alcohol) $\xrightarrow[\text{(ii) Et}_2\text{SO}_4]{\text{(i) Base}}$ $\xrightarrow[h\upsilon]{\text{Air}}$ $\xrightarrow[\text{KCNS}]{\text{FeSO}_4}$ Red colour}
\opt{D}{\qfigi{p1_chem_s2_q12_D}}
\end{qoptions}
\sect{SECTION-II}
\qn{1} A saturated solution in AgA (K$_\text{sp}$ = 3 $\times$10$^{-14}$) and AgB (K$_\text{sp}$ = 1 $\times$ 10$^{-14}$) has conductivity of 375 $\times$ 10$^{-10}$ Scm$^{-1}$ and limiting molar conductivity of Ag$^+$ and A$^-$ are 60 Scm$^2$ mol$^{-1}$ and 80 Scm$^2$ mol$^{-1}$ respectively then what will be the limiting molar conductivity of B$^-$ (in Scm$^2$ mol$^{-1}$)
\qn{2} Find the number of species in which number of anions satisfying only primary valency is / are lesser than anion which satisfies secondary valency as well as primary valency.\par
Cr.6NH$_3$.Cl$_3$ \quad Cr.5NH$_3$.Cl$_3$ \quad Cr.3NH$_3$.Cl$_3$\par
Pt.6NH$_3$.Cl$_4$ \quad Pt.5NH$_3$.Cl$_4$ \quad Pt.4NH$_3$.Cl$_4$\par
Pt.2NH$_3$.Cl$_4$
\qn{3} \qfigi{p1_chem_s3_q03_1}\par
How many optically active compounds having molecular formula C$_8$H$_8$Br$_8$ will be obtained in above reaction ?
\qn{4} In cubical close packing (ccp) arrangement of solid what will be total effective number of voids present in the given plane in the given diagram.\par
\qfig{p1_chem_s3_q04_1}
\qn{5} How many of the following complex do not exhibit optical activity but having diamagnetic characters.\par
(a) Cis--[Co(en)$_2$(H$_2$O)$_2$]$^{3+}$ \quad (b) Cis--[Pt(NH$_3$)$_2$Cl$_2$]\par
(c) Trans--[Co(en)$_2$(H$_2$O)$_2$]$^{3+}$ \quad (d) [Ni(CO)$_4$]\par
(e) cis[Ni(NH$_3$)$_4$(H$_2$O)$_2$]SO$_4$\par
(f) [Zn(Gly)$_2$]$^\text{O}$
\qn{6} Consider the following compound :\par
\qfigi{p1_chem_s3_q06_2}\qfigi{p1_chem_s3_q06_4}\qfigi{p1_chem_s3_q06_3}\qfigi{p1_chem_s3_q06_1}\par
\qfigi{p1_chem_s3_q06_7}\qfigi{p1_chem_s3_q06_5}\qfigi{p1_chem_s3_q06_6}\qfigi{p1_chem_s3_q06_8}\par
Serial number of most basic compound is (P)\par
Serial number of most acidic compound is (Q)\par
Total number of aromatic compound = (R)\par
Number of compound is represented in bracket.\par
Find the value of P + Q + R.
\partsec{3}{MATHEMATICS}
\sect{SECTION-I (i)}
\qn{1} The value of $\tan^{-1}\left(\dfrac{1-\sin5}{\cos5}\right)$ is equal to :
\begin{qoptions}
\opt{A}{$\dfrac{5}{2}-\dfrac{\pi}{4}$}
\opt{B}{$\dfrac{\pi}{4}-\dfrac{5}{2}$}
\opt{C}{$\dfrac{5}{2}-\dfrac{5\pi}{4}$}
\opt{D}{$\dfrac{5\pi}{4}-\dfrac{5}{2}$}
\end{qoptions}
\qn{2} Let for a sequence a$_\text{n}$ = $\dfrac{1}{\sqrt{\text{n}}}$, n $\in$ N then $\left[\displaystyle\sum_{\text{n}=1}^{80}\text{a}_\text{n}\right]$ is equal to (where [.] represent greatest integer function) :
\begin{qoptions}
\opt{A}{16}
\opt{B}{17}
\opt{C}{18}
\opt{D}{15}
\end{qoptions}
\qn{3} Minimum value of expression $|\sin\text{x}+\cos\text{x}+\tan\text{x}+\cot\text{x}+\sec\text{x}+\text{cosec x}|$ is :
\begin{qoptions}
\opt{A}{$2\sqrt{2}-1$}
\opt{B}{$2\sqrt{2}+1$}
\opt{C}{$3\sqrt{2}+2$}
\opt{D}{2}
\end{qoptions}
\qn{4} If $\displaystyle\int_0^{2\pi}\dfrac{|\sin\theta|\,\text{d}\theta}{1+|\cos2\theta|}$ = 2$\left(\sqrt{2}-1+\ell\text{na}\right)$, then (a -- 1)$^2$ is
\begin{qoptions}
\opt{A}{0}
\opt{B}{2}
\opt{C}{5}
\opt{D}{8}
\end{qoptions}
\qn{5} Number of points of discontinuity of the function $f$(x) = sgn(x$^2$ -- x + 1) + [x$^2$ -- x + 1] in x $\in$ [0, 3] ([.] denotes greatest integer function, sgn \{.\} denotes signum function) -
\begin{qoptions}
\opt{A}{7}
\opt{B}{6}
\opt{C}{14}
\opt{D}{8}
\end{qoptions}
\qn{6} If probability that 2$^3$3$^5$5$^4$7$^2$11$^\alpha$ be resolved as product of two factors such that their highest common factor is 1 or 5 or 7 is $\left(\dfrac{1}{27}\right)$, then $\alpha$ is
\begin{qoptions}
\opt{A}{1}
\opt{B}{3}
\opt{C}{5}
\opt{D}{8}
\end{qoptions}
\sect{SECTION-I (ii)}
\qn{7} A parabola touches x-axis and y-axis at (1, 0) and (0, 1) respectively then :
\begin{qoptions}
\opt{A}{Equation of parabola is (y $-$ x)$^2$ = 2x + 2y $-$ 1}
\opt{B}{Coordinates of focus parabola is $\left(\dfrac{1}{2},\dfrac{1}{2}\right)$}
\opt{C}{Equation of directrix of parabola is x + y = 0}
\opt{D}{Length of latus rectum of parabola is 2}
\end{qoptions}
\qn{8} Let I$_\text{n}$ = $\displaystyle\int_{-\pi}^{\pi}\dfrac{1}{(1+2^\text{x})}\left(\dfrac{\sin\frac{\text{nx}}{2}}{\sin\frac{\text{x}}{2}}\right)^2$dx for n $\in$ N then :
\begin{qoptions}
\opt{A}{I$_{\text{n}+1}$ $-$ I$_\text{n}$ = $\pi$, $\forall$ n $\in$ N}
\opt{B}{(I$_0$)$^2$ $-$ (I$_1$)$^2$ + (I$_2$)$^2$ $-$ (I$_3$)$^2$ + ....+(I$_{100}$)$^2$ = 5051 $\pi^2$}
\opt{C}{$\displaystyle\sum_{\text{m}=1}^{9}$ I$_{2\text{m}}$ = 90 $\pi$}
\opt{D}{I$_1$ $-$ I$_2$ + I$_3$ $-$ I$_4$ + ....$-$I$_{100}$ = 50 $\pi$}
\end{qoptions}
\qn{9} S = (1)$^{\frac{1}{3}}$ + (2)$^{\frac{1}{3}}$ + ..... + (1000)$^{\frac{1}{3}}$ then
\begin{qoptions}
\opt{A}{S $>$ 7500}
\opt{B}{S $<$ 7510}
\opt{C}{S $<$ 5000}
\opt{D}{S $>$ 5010}
\end{qoptions}
\qn{10} Which of the following is/are correct :
\begin{qoptions}
\opt{A}{The total number of solutions of the equation xyz = 10080, where x . y, z $\in$ Z is 3600}
\opt{B}{The total number of solutions of the equation xyz = 10080, where x. y, z $\in$ N is 900}
\opt{C}{Number of zeroes at the end of $^{28}$C$_{16}$ is 1}
\opt{D}{Number of zeroes at the end of $^{28}$C$_{16}$ is 0}
\end{qoptions}
\qn{11} Let f(x) = 2x$^3$ -- 3x$^2$ -- 12x + 8 then subsets of solution set of the inequality 2(f(x))$^3$ -- 3(f(x))$^2$ -- 12f(x) + 8) $>$ f(2x -- 7) \ is :
\begin{qoptions}
\opt{A}{(1, 2)}
\opt{B}{$\left(-\dfrac{5}{2},-1\right)$}
\opt{C}{(2, $\infty$)}
\opt{D}{(3, $\infty$)}
\end{qoptions}
\qn{12} If z$^2$ + $\bar{\text{z}}$ -- z = 4i and z does not lie in third quadrant, then-
\begin{qoptions}
\opt{A}{$|$z$|$ = $\sqrt{2}$}
\opt{B}{$|$z$|$ = 2$\sqrt{2}$}
\opt{C}{arg(z) = $\dfrac{\pi}{4}$}
\opt{D}{arg(z) = $\dfrac{3\pi}{4}$}
\end{qoptions}
\sect{SECTION-II}
\qn{1} Let A be the set of all integral solution of the inequality x$^2$ -- $|$x$|$ -- 2 $\leq$ 0 and $\vec{\text{v}}_1$ = $\hat{\text{i}}$ + $\hat{\text{j}}$ + $\hat{\text{k}}$, $\vec{\text{v}}_2$ = a$\hat{\text{i}}$ + b$\hat{\text{j}}$ + c$\hat{\text{k}}$ (where a, b, c $\in$ A) be two non zero vectors such that $\vec{\text{v}}_2$ is orthogonal to $\vec{\text{v}}_1$ then number of such possible vector $\vec{\text{v}}_2$ is :
\qn{2} Number of solution of the equation $\left\{\sin^{-1}(\sin\text{x})\right\}$ = $\dfrac{\text{x}}{10}$, in [0, 2$\pi$) is, where \{.\} denotes fractional part function :
\qn{3} Coefficient of x$^{99}$ in (x $-$ 1) (x$^2$ $-$ 2) ..... (x$^{14}$ $-$ 14) is :
\qn{4} Find the least positive integral value of `m' such that $\left({}^{2\text{n}}\text{C}_\text{n}\right)^{1/\text{n}}$ $<$ m \ $\forall$ n $\in$ N :
\qn{5} Let `$\alpha$' and `$\beta$' are the roots of quadratic x$^2$ $-$ 5x $-$ 2 = 0 then the value of $\dfrac{\alpha^{2023}-8\alpha^{2017}+\beta^{2023}-8\beta^{2017}}{\alpha^{2020}+\beta^{2020}}$ is equal to :
\qn{6} If A and B are non-singular matrices of order three such that adj(AB) = $\begin{bmatrix}1&1&1\\1&\text{p}&1\\1&1&\text{p}\end{bmatrix}$ and $|\text{B}^2\text{adj(A)}|$ = p$^2$ $-$ 3, then p is
\raggedcolumns
\papersec{2}
\partsec{1}{PHYSICS}
\sect{SECTION-I}
\qn{1} Radii of the inner and outer concentric conducting spheres of a spherical capacitor are a and b. One half of the space between the spheres is filled with a linear isotropic dielectric of permittivity $\varepsilon_1$ and the other half with another linear isotropic dielectric of permittivity $\varepsilon_2$ as shown in the figure. The outer sphere is given a total charge +Q and the inner sphere a total charge $-$Q. Which of the following statements are correct?\par
\qfig{p2_phy_s1_q01_1}
\begin{qoptions}
\opt{A}{Electric fields in both the dielectrics at the same radial distance are not equal.}
\opt{B}{Electric field intensities in both the dielectrics at radial distance r from center are given by $\text{E}_1 = \text{E}_2 = \dfrac{\text{Q}}{2\pi(\varepsilon_1+\varepsilon_2)\text{r}^2}$}
\opt{C}{Potential difference between the spheres is given by V = $\dfrac{\text{Q(b}-\text{a)}}{2\pi(\varepsilon_1+\varepsilon_2)\text{ab}}$}
\opt{D}{Charge densities on both the halves of a sphere are equal.}
\end{qoptions}
\qn{2} Figure shows an alternating current circuit. Voltage of source is given by V = 50sin($\omega$t + $\delta$). Mark the CORRECT statement(s) :-\par
\qfig{p2_phy_s1_q02_1}
\begin{qoptions}
\opt{A}{Phase of current in B leads voltage of A by a phase $\theta$ then $\theta$ should be $\theta$ = 0 rad.}
\opt{B}{Phase of current in D leads voltage of C by phase $\theta$ then should be $\theta = \dfrac{+\pi}{2}$ rad.}
\opt{C}{Phase of voltage of D leads current of B by phase $\theta$ then $\theta$ should be $\dfrac{-\pi}{2} < \theta < 0$ rad.}
\opt{D}{Phase of current of C leads current of A by phase $\theta$ then $\theta$ should be $0 < \theta < \dfrac{\pi}{2}$ rad.}
\end{qoptions}
\qn{3} A uniform solid cylinder of mass M and radius R rolls down a rough inclined plane as shown in figure. As the cylider rolls it winds up a light string which passes over a light pulley and connected to mass m. The tension in the string will be :\par
\qfig{p2_phy_s1_q03_1}
\begin{qoptions}
\opt{A}{$\dfrac{(3+4\sin\theta)\text{Mmg}}{3\text{M}+8\text{m}}$}
\opt{B}{$\dfrac{(3+5\sin\theta)\text{Mmg}}{3\text{M}+8\text{m}}$}
\opt{C}{$\dfrac{(3\text{cosec}\,\theta+4)\text{Mmg}}{(3\text{M}+8\text{m})\text{cosec}\,\theta}$}
\opt{D}{$\dfrac{(3\text{cosec}\,\theta+5)\text{Mmg}}{(3\text{M}+8\text{m})\text{cosec}\,\theta}$}
\end{qoptions}
\qn{4} Figure shows two reservoirs containing two liquids of masses 4m and 2m and their specific heat capacities are S and 2S respectively. Their initial temperatures are 4T$_0$ and T$_0$ respectively. The containers are joined by a conducting rod of thermal conductivity K, length $\ell$ and cross-section area A. Specific heat capacity of the rod is negligible. (Neglect radiation)\par
\qfig{p2_phy_s1_q04_1}
\begin{qoptions}
\opt{A}{The temperature of mid-point `P' of the rod is always constant.}
\opt{B}{Heat extracted from container 1 = Heat supplied to container 2.}
\opt{C}{After time t$_0$ = $\dfrac{2\text{ms}\ell}{\text{kA}}\ell$n2 the temperature difference between containers becomes half of the initial temperature difference.}
\opt{D}{After time t$_0$ = $\dfrac{\text{ms}\ell}{\text{kA}}\ell$n2 the temperature difference between containers becomes half of the initial temperature difference.}
\end{qoptions}
\qn{5} Shown in figure there is a simple network of three resistors R$_1$ = 1$\Omega$, R$_2$ = 2$\Omega$ and R$_3$ = 3$\Omega$. A battery of emf 0.3 volt is connected. The dimension of the loop are given. In loop (1) inward magnetic field B$_1$ = e$^{-2t}$ is applied and in loop (2) outward magnetic field B$_2$ = 2t is applied. Find the initial current flown through the resistor R$_3$ and R$_2$. (t is time)\par
\qfig{p2_phy_s1_q05_1}
\begin{qoptions}
\opt{A}{initial current in R$_3$ = $\dfrac{9}{110}$A}
\opt{B}{initial current in R$_2$ = $\dfrac{6}{110}$A}
\opt{C}{initial current in R$_2$ = $\dfrac{8}{110}$A}
\opt{D}{initial current in R$_3$ = $\left(\dfrac{1}{110}+\dfrac{1}{20}\right)$A}
\end{qoptions}
\qn{6} A particle is dropped along the axis from a height f/2 on a concave mirror of focal length f as shown in figure. The acceleration due to gravity is g. The maximum speed of image is :-\par
\qfig{p2_phy_s1_q06_1}
\begin{qoptions}
\opt{A}{$\dfrac{3}{8}\sqrt{12\text{fg}}$}
\opt{B}{$\dfrac{3}{4}\sqrt{5\text{fg}}$}
\opt{C}{$\dfrac{3}{4}\sqrt{3\text{fg}}$}
\opt{D}{$\dfrac{3}{16}\sqrt{7\text{fg}}$}
\end{qoptions}
\sect{SECTION-II}
\qn{1} A cylinder of mass M and raidus R moves with constant speed v through a region of space that contains dust particles of mass m which is at rest. There are n number of particles per unit volume. The cylinder moves in a direction perpendicular to its axis. Assume m $\ll$ M, and assume the particles do not interact with each other. All the collision takes place is perfectly elastic and the surface of the cylinder is smooth. The drag force per unit length of the cylinder require to maintain the speed v constant for the cylinder when it has entered in the region is $\dfrac{\text{K}}{3}$nm Rv$^2$. Find the value of K.\par
\qfig{p2_phy_s2_q01_1}
\qn{2} Half life of a radioactive substance is 2.34 minute. It is produced at a constant rate of 10$^8$ nuclei per sec. How soon after (in minute), the beginning of production will its activity be equal to 7.5 $\times$ 10$^7$ disintegration per sec ?
\qn{3} A point charge q is kept at the center of circular plane surface of cone. A shaded section ABDCA shown at the slant surface of cone. If electric flux through the shaded section is $\dfrac{\text{n}}{\text{m}} \times \dfrac{\text{q}}{80\pi\varepsilon_0}$ where n and m minimum possible integer, find (n + m). [length AB = $\dfrac{\text{h}}{2}$]\par
\qfig{p2_phy_s2_q03_1}
\qn{4} The pitch of screw gauge is 0.75 mm \& there are 50 divisions on the circular scale. When nothing is put between the studs it shows the following where zero of mains scale is not visible. Now thickness of a metal plate is measured and when the plate is put between the studs it shows the following. The thickness of metal plate is (in mm) :-\par
\qfig{p2_phy_s2_q04_1}
\qn{5} A resistance (R), inductance (L) and capacitance (C) are connected in series to an ac source of voltage v having variable frequency. The energy delivered by the source to the circuit during one period if the operating frequency is twice the resonance frequency is $\dfrac{\pi\text{R}\sqrt{\text{LC}}\,\text{V}^2}{\text{R}^2+\text{P}\left(\frac{\text{L}}{\text{C}}\right)}$, find P.
\qn{6} A binary star is a system of two stars moving around the centre of inertia of the system due to mutual gravitational force. If the total mass of two stars = M$_0$ and the period of revolution is T, find the distance in km between stars.(Take GM$_0$ = 8 $\times$ 10$^6$ and T = 2$\pi$ $\times$ 10$^3$ seconds).
\sect{SECTION-III}
\qn{1} In the circuit shown in the figure, S$_1$ remains closed for a long time and S$_2$ remains open. Now S$_2$ is closed and S$_1$ is opened. Find out the $\left|\dfrac{\text{di}}{\text{dt}}\right|$ (in Amp/sec) just after that moment, where i is the current through the inductance.\par
\qfig{p2_phy_s3_q01_1}
\qn{2} In a series LCR circuit, half power frequencies (frequency at which power disspited is half the maximum power disspited in the circuit)are found to be 90 rad/s and 40 rad/s. Find the value of $\dfrac{\omega_0}{10}$, where $\omega_0$ is the resonance frequency.
\qn{3} Water at temperature 40\textdegree C flows from a tap T into a heated container C. The container has a heating element (a resistor R) which is generating heat at the rate of P, that may be varied. The rate of water inflow from tap is m = $\dfrac{1000}{7}\ell$/min. The heated generated is sufficient so that the water in the container is boiling and getting converted into steam at a steady rate. The minimum power P (in MW) that must be generated as heat in the steady state in resistor R so that the amount of liquid water in the container neither increases nor decreases with time ? (Neglect other losses of heat, such as conduction from the container to the air and heat capacity of container)\par
For water, specific heat c = 4.2 kJ kg$^{-1}$ K$^{-1}$, latent heat of vaporization L$_\text{vap}$ = 2.268 MJ kg$^{-1}$, density $\rho$ = 1000 kg m$^{-3}$.\par
\qfig{p2_phy_s3_q03_1}
\qn{4} Let the radius of ground state H-atom (as per Bohr's atomic model) be r$_0$ and its ionisation energy be E$_0$. A hydrogen like system of e$^+$ $-$ e$^-$ (electron-positron) is called positronium. If the radius of ground state positronium is r and ionisation energy is E, then find the value of $\dfrac{\text{rE}_0}{\text{r}_0\text{E}}$.
\qn{5} A concave mirror of radius of curvature 50 cm lies at origin of the coordinate system such that its axis coincides with x axis. A small straight object lies on the x-axis along the line 5y + 4x $-$ 25 = 0, where x and y are in cm. Assuming mirror formula to be valid for the image formation of this object, calculate the value of 5m where m is slope of the line along which the image lies.
\qn{6} A boy runs on a circular path of radius R = 28 m with a constant speed u = 4 m/s. Another boy starts from the centre of the path to catch the first boy. The second boy always remains on the radius connecting the centre of the circle and the first boy and maintains magnitude of his velocity constant V = 4 m/s. If the time of chase is (10 + x) sec then x =.....; $\left(\pi = \dfrac{22}{7}\right)$
\raggedcolumns
\partsec{2}{CHEMISTRY}
\sect{SECTION-I}
\qn{1} \qfig{p2_chem_s1_q01_1}
$\Delta\text{H}^0_\text{f}$ for P = 36 KJ/mole\par
$\Delta\text{H}^0_\text{f}$ of monomeric form of P = 100 KJ/mole\par
$\Delta\text{H}^0_\text{atom}$ for Q(g) = 436 KJ/mole\par
$\Delta\text{H}^0_\text{atom}$ for S(solid) = 565 KJ/mole\par
Then :
\begin{qoptions}
\opt{A}{Bond enthalpy of each of 3c$-$2e bond in `P' is 32 KJ/mole.}
\opt{B}{$\text{P} + \text{NH}_3 \xrightarrow[400\text{K}]{1:2\text{ ratio}} \text{T} + \text{Q}$, here `T' is resonance stabilised.}
\opt{C}{Bond enthalpy of each 3c$-$2e bond in `P' is 455 KJ/mole.}
\opt{D}{In the option (B) if NH$_3$ is in excess at high temperature, product is having graphite like structure.}
\end{qoptions}
\qn{2} If Compounds $\text{HO}-\overset{\text{O}}{\overset{\|}{\text{C}}}-(\text{CH}_2)_\text{n}-\overset{\text{O}}{\overset{\|}{\text{C}}}-\text{OH}$ ; where n = 0 is (P), n = 1 is (Q) \& n = 2 is (R). The least acidic compound is taken for following sequence of reactions.\par
\qfig{p2_chem_s1_q02_1}
\begin{qoptions}
\opt{A}{`T' is more basic than `S'}
\opt{B}{`U' is lactam.}
\opt{C}{`T' is less basic than `S'}
\opt{D}{In the reaction S $\xrightarrow[\text{BF}_3]{\text{Ac}_2\text{O}}$ product, the substitution at position 3 takes place in the compound `S'.}
\end{qoptions}
\qn{3} Which of the following is correct when an ideal gas having initial pressure P, volume V and temperature T is allowed to expand adiabatically until its volume becomes 5.66 V while its temperature falls to half ? (Given : 10$^{0.7588}$ = 5.66, log 2 = 0.301)
\begin{qoptions}
\opt{A}{The gas is diatomic}
\opt{B}{The gas is monoatomic}
\opt{C}{The work done by the gas is 1.25 PV}
\opt{D}{All are correct}
\end{qoptions}
\qn{4} Which statements is/are correct among the following ?
\begin{qoptions}
\opt{A}{BF$_3$ is strongest Lewis acid among BF$_3$, BCl$_3$, BBr$_3$.}
\opt{B}{Between NH$_3$ and PH$_3$, NH$_3$ is better electron donor because NH$_3$ has more directional lone pair.}
\opt{C}{In P$_4$O$_{10}$ there are 5, P$-$O$-$P bonds.}
\opt{D}{ClOF$_3$ has same shape as XeO$_2$F$_2$.}
\end{qoptions}
\qn{5} Which of the following will be soluble in NaHCO$_3$ solution and more acidic than \qfigi{p2_chem_s1_q05_1} ?
\begin{qoptions}
\opt{A}{\qfigi{p2_chem_s1_q05_A}}
\opt{B}{\qfigi{p2_chem_s1_q05_B}}
\opt{C}{\qfigi{p2_chem_s1_q05_C}}
\opt{D}{\qfigi{p2_chem_s1_q05_D}}
\end{qoptions}
\qn{6} In order to study effect of temperature on rate of reaction, another term ``$\sigma$'' is defined as $\sigma = \dfrac{1}{\text{K}}\dfrac{\text{dk}}{\text{dT}}$ [where k is rate constant and T represents temperature]. If for a reaction value of $\sigma = \dfrac{1.25\times10^6}{\text{T}^3}\ \text{K}^{-1}$ [T should be in Kelvin], then identify the correct statement(s).
\begin{qoptions}
\opt{A}{Activation energy for the given reaction at 250 Kelvin is 10 Kcal/mol.}
\opt{B}{Activation energy for the given reaction at 2000 Kelvin is 1.25 Kcal/mol.}
\opt{C}{Rate of increase of rate constant with temperature is higher at lower temperature range than at higher temperature range.}
\opt{D}{The value of $\dfrac{\text{dK}}{\text{dT}}$ at 2000 Kelvin is 0.625 $\times$ K.}
\end{qoptions}
\sect{SECTION-II}
\qn{1} Give the maximum \% of R$-$2$-$bromopentane that could be formed from the reaction of bromine with n$-$pentane (Assuming reactivity order of 3\textdegree H : 2\textdegree H : 1\textdegree H = 1600 : 82 : 1)
\qn{2} $\text{N}_2\text{(g)} + 3\text{H}_2\text{(g)} \rightleftharpoons 2\text{NH}_3\text{(g)}$ \quad $\Delta\text{G}^0_{\text{f(NH}_3)}$ = $-$15.53 kJ/mole at 300 K.\par
If after 1 hour of mixing reactants, partial pressureof N$_2$ = 50 Bar, H$_2$ = 2 Bar, NH$_3$ = 200 Bar \& $\Delta\text{H}_\text{reaction}$ is $-$4.2 kJ, find $\Delta$S of reaction (in JK$^{-1}$) (T = 300 K = constant).\par
($\ell$n10 = 2.3, R = $\dfrac{25}{3}$JK$^{-1}$ mol$^{-1}$).
\qn{3} Number of Cr-O linkages in chromate ion = x\par
Number of Cr-O linkages in dichromate ion = y\par
Number of Cr-O linkages in chromyl chloride = z\par
Determine, (x + y + z).
\qn{4} Halide of Phosphorus PX$_\text{a}$ is found to exist in ionic form in solid state, undergoes complete hydrolysis to give two acids H$_\text{b}$PO$_\text{c}$ and HX. Salt NaX on reaction with conc H$_2$SO$_4$ liberates gas X$_2$. Elemental form of X$_2$ at room temperature is liquid. Value of $\left(\dfrac{\text{at. wt. of X (in amu)}}{\text{a}+2\text{b}+3\text{c}}\right)$
\qn{5} An ideal mixture of liquids A and B with 2 moles of A and 2 moles of B has a total vapour pressure of 1 atm at a certain temperature. Another mixture with 1 mole of A and 3 moles of B has a vapour pressure greater than 1 atm. But if 4 moles of C are added to the second mixture, the vapour pressure comes down to 1 atm. Vapour pressure of C in pure state is 0.8 atm. Calculate the vapour pressure of pure B (in atm):-
\qn{6} Identify total number of $\alpha$-`H' on major product of following reaction?\par
\qfig{p2_chem_s2_q06_1}
\sect{SECTION-III}
\qn{1} The emf of the following cell is given as E$_\text{Cell}$ = $-$a $\times$ 10$^{-1}$ V\par
Pt $|$ H$_2$(1 atm) $|$ CH$_3$COOH (0.1 M) $\|$ NH$_4$OH(0.1 M) $|$ H$_2$(1 atm) $|$ Pt\par
K$_\text{a}$ for CH$_3$COOH = 10$^{-5}$; K$_\text{b}$ for NH$_4$OH = 10$^{-5}$, Given $\dfrac{2.303\,\text{RT}}{\text{F}}$ = 0.0625\par
What is the value of `a' ?
\qn{2} Identify how many ``Ore names'', ``Type of Anion'', ``Metal(s) present in mineral'' are correctly related\par
{\small\setlength{\tabcolsep}{3pt}\begin{tabular}{@{}l l l l@{}}
(i) & Magnetite & oxide & Fe\\
(ii) & Gibbsite & hydoxide & Al\\
(iii) & Copper Pyrite & Sulphide & Fe\\
(iv) & Galena & sulphide & Hg\\
(v) & Cerussite & Sulphate & Pb\\
(vi) & Lunar Caustic & hydroxide & Ag\\
(vii) & Dolomite & carbonate & Ca\\
(viii) & Dolomite & hydroxide & Mg\\
\end{tabular}}
\qn{3} Number of reactions in which aldehyde can be obtained as one of the product.\par
(i) $\text{CH}_3\text{--CH}_2\text{OH} \xrightarrow[300\text{C}]{\text{Cu}}$\par
(ii) $\text{CH}_3\text{--CH}_2\text{--CH}_2\text{--OH} \xrightarrow[\text{Pyridine}]{\text{CrO}_3}$\par
\qfig{p2_chem_s3_q03_2}
\qn{4} For an ideal gas sample (fixed mole taken in fixed volume container) collision frequency (Z) \& mean free path ($\lambda$).\par
(i) Both Z \& $\lambda$ change with change in T.\par
(ii) Z changes with change in T.\par
(iii) $\lambda$ changes with change in T.\par
(iv) Both are independent of temperature T.\par
(v) Heavier gas (more mole weight) exibits high Z compared to lighter gas, at some temperature.\par
(vi) Heavier gas exibits more $\lambda$ compared to lighter gas at same temperature.\par
(vii) Molecular weight of gas does not have any effect of value of $\lambda$.\par
Identify number of correct statements from above ?
\qn{5} Find the total number of optically active compounds\par
\qfigi{p2_chem_s3_q05_2} \quad $\text{CH}_3-\text{CH}=\text{C}=\text{CH}-\text{CH}_3$, \quad \qfigi{p2_chem_s3_q05_1}\par
\qfigi{p2_chem_s3_q05_3}
\qn{6} How many of the following statements are correct ? (R = atomic radius)\par
(1) In CCP, distance between two nearest octahedral voids is 2R\par
(2) In CCP, distance between nearest tethrahedral and octahedral voids is $2\sqrt{3}$R\par
(3) In CCP, distance between two nearest voids is $\sqrt{2}$R\par
(4) In CCP, distance between next nearest tetrahedral void is $\sqrt{5}$R\par
(5) In HCP, on each edge, centres of two tetrahedral voids and one octahedral void is located\par
(6) In HCP, if a plane is passed parallel to top face but at $\dfrac{3\text{h}}{4}$ from bottom, it pass through centres of octahedral voids. (h is the height of prismatic unit cell)\par
(7) In HCP, top plane pass through centres of both tetrahedral and octahedral voids.
\raggedcolumns
\partsec{3}{MATHEMATICS}
\sect{SECTION-I}
\qn{1} Which of the following is/are correct ?
\begin{qoptions}
\opt{A}{$\displaystyle\int_1^3 \dfrac{[\text{x}^2]}{[\text{x}^2-8\text{x}+16]+[\text{x}^2]}\,\text{dx} = 1$}
\opt{B}{$\displaystyle\int_1^3 \dfrac{[\text{x}^2]}{[\text{x}^2-8\text{x}+16]+[\text{x}^2]}\,\text{dx} = 2$}
\opt{C}{$\displaystyle\int_0^2 \left([\text{x}]-\left[\dfrac{\text{x}}{2}\right]\right)\text{dx} = 1$}
\opt{D}{$\displaystyle\int_0^2 \left([\text{x}]-\left[\dfrac{\text{x}}{2}\right]\right)\text{dx} = 2$ (where [.] denotes greatest integer function)}
\end{qoptions}
\qn{2} Let the function y = $f$(x) satisfies the differential equaiton xdy = e$^\text{x}$(x $\ell$n x + 1)dx and passes through (1, 0), then :
\begin{qoptions}
\opt{A}{$f$(x) $>$ 0 $\forall$ x $\in$ (1, $\infty$)}
\opt{B}{$\displaystyle\int \text{e}^{-\text{x}} f(\text{x})\,\text{dx} = \text{x}\ell\text{nx} - \text{x} + \text{c}$ (where `c' is constant of integration)}
\opt{C}{$\displaystyle\lim_{\text{x}\to1}\dfrac{f(\text{x})}{\text{x}-1} = \text{e}$}
\opt{D}{$\displaystyle\lim_{\text{x}\to1}(1+f(\text{x}))^{\frac{1}{\text{x}-1}} = \text{e}^\text{e}$}
\end{qoptions}
\qn{3} If the position vector of a point P is $\vec{\text{r}} = \text{x}\hat{\text{i}}+\text{y}\hat{\text{j}}+\text{z}\hat{\text{k}}$ where x, y, z $\in$ N and projection of $\vec{\text{r}}$ on $\vec{\text{a}} = \hat{\text{i}}+\hat{\text{j}}+\hat{\text{k}}$ is $\dfrac{10}{\sqrt{3}}$, then number of possible positions of P is also equal to :
\begin{qoptions}
\opt{A}{Number of natural solutions of x + y + z = 7.}
\opt{B}{Total number of outcomes when a pair of dice are rolled once.}
\opt{C}{Number of divisors of 1800.}
\opt{D}{Number of cyphers at the end of (120)!}
\end{qoptions}
\qn{4} Let $f$(x) = $\begin{cases}3^\text{x} & ;\ \text{x}\leq0\\ (\text{x}+\lambda) & ;\ 0<\text{x}<2\\ 3+(\text{x}-2)^2 & ;\ \text{x}\geq2\end{cases}$ is continuous for all x $\in$ R, then which of the following is true -
\begin{qoptions}
\opt{A}{Range of $f$(x) is (0, $\infty$)}
\opt{B}{$f$(x) = 1 has one solution}
\opt{C}{$f$(x) = $|f$(x)$|$ for all x $\in$ R}
\opt{D}{$f$(x) = 1 has two solutions}
\end{qoptions}
\qn{5} Let $f$ : R $\to$ R be a continuous function satisfying x$^2$ $-$ x($f$(x) + 5) + 6 + 2$f$(x) = 0 $\forall$ x $\in$ R. Then which of the following holds good ?
\begin{qoptions}
\opt{A}{$f$ is many-one function}
\opt{B}{$f$(8) = 5}
\opt{C}{Area bounded by y = $f$(x) and co-ordinates axes is $\dfrac{9}{2}$}
\opt{D}{$f$(x) = $f^{-1}$(x) is satisfied for exactly one value of x.}
\end{qoptions}
\qn{6} If f(x) = e$^{-\text{ax}}\left(\text{x}^2+\text{ax}+\text{b}\right)$ (a $>$ 0) has points of extremum at points x = $-$2 \& x = 1, then
\begin{qoptions}
\opt{A}{a = 2, b = $-$1}
\opt{B}{a = 1, b = $-$2}
\opt{C}{$f''(-2)<0$ \& $f''(1)>0$}
\opt{D}{$f''(-2)>0$ \& $f''(1)<0$}
\end{qoptions}
\sect{SECTION-II}
\qn{1} Let $f$(x) = $\begin{cases}\dfrac{\text{e}^{-\text{x}^2/2}-\cos\text{x}}{\text{x}\,\ell\text{n}(1+\text{x})\sin\text{x}(\text{e}^\text{x}-1)} & ;\ \text{x}\neq0\\[1.5ex] \text{k} & ;\ \text{x}=0\end{cases}$ be continuous at x = 0, then k equals -
\qn{2} The probability that Ramu's Gardner will forget to water the rose bush is $\dfrac{2}{3}$. The rose bush is in questionable condition. Any how if watered, the probability of its withering is $\dfrac{1}{2}$ and if not watered then the probability of its withering is $\dfrac{3}{4}$. Ramu went out of station and after returning he finds that rose bush has withered. The probability that the Gardner did not water the rose bush is-
\qn{3} Let L$_1$ : $\dfrac{\text{x}-1}{2}=\dfrac{\text{y}-2}{-1}=\dfrac{\text{z}-3}{3}$, L$_2$ : $\dfrac{\text{x}-1}{-1}=\dfrac{\text{y}-2}{3}=\dfrac{3(\text{z}-3)}{5}$ and L$_3$ : $\dfrac{\text{x}-1}{-32}=\dfrac{\text{y}-2}{-19}=\dfrac{\text{z}-3}{15}$ be three lines. A plane intersecting these lines at A, B and C respectively such that PA = 2, PB = 3 and PC = 6, where P $\equiv$ (1, 2, 3), then the perpendicular distance `d' of the plane from the point P is -
\qn{4} Let $f$(x) = $\dfrac{1}{\text{e}^\text{x}+8\text{e}^{-\text{x}}+4\text{e}^{-3\text{x}}}$, g(x) = $\dfrac{1}{\text{e}^{3\text{x}}+8\text{e}^\text{x}+4\text{e}^{-\text{x}}}$ and $\displaystyle\int(f(\text{x})-2\text{g}(\text{x}))\,\text{dx} = \text{h}(\text{x})+\text{c}$ (where `c' is constant of integration) and $\displaystyle\lim_{\text{x}\to\infty}\text{h}(\text{x}) = \dfrac{\pi}{4}$. If h(0) = $\dfrac{1}{\text{a}}\tan^{-1}\left(\dfrac{\text{b}}{\text{c}}\right)$ (where a $\in$ N, b and c are coprime), then the value of $\left(\dfrac{\text{a}+\text{c}}{\text{b}}\right)$ is:
\qn{5} If x and y are the solutions of the equation 12sin x + 5cos x = 2y$^2$ $-$ 8y + 21, the least positive value of $\dfrac{\text{x}}{\text{y}}$ is $\dfrac{1}{\text{a}}\cot^{-1}\left(\dfrac{\text{b}}{\text{c}}\right)$ then the value of $\left(\dfrac{\text{ab}}{\text{c}}\right)$ is :
\qn{6} Let the line y $-$ $\sqrt{3}$x + 3 = 0 cuts the parabola 2y$^2$ = 2x + 3 at A and B. If P($\sqrt{3}$, 0), then value of $|$PA $-$ PB$|$ is [where PA denotes distance between points P and A] :
\sect{SECTION-III}
\qn{1} Suppose $\alpha,\beta,\gamma$ are complex numbers satisfying the following system of the equations $\alpha+\beta+\gamma = 6$, $\alpha^3+\beta^3+\gamma^3 = 87$ and $(\alpha+1)(\beta+1)(\gamma+1) = 33$.\par
If $\dfrac{1}{\alpha}+\dfrac{1}{\beta}+\dfrac{1}{\gamma} = \dfrac{\text{m}}{\text{n}}$ (where m \& n $\in$ N are relatively prime), then the value of $\left[\dfrac{\text{n}}{\text{m}}\right]$ is\par
(where [.] denotes greatest integer function)
\qn{2} Let z$_\text{n}$ denotes the n$^\text{th}$ term of a squence of complex numbers defined by z$_\text{n}$ = (5 + 12i)$\left(\dfrac{\text{i}}{2}\right)^\text{n}$. Let S be the intersection of the closed unit disk centered at origin and the fourth quadrant in the complex plane, then the smallest value of `n' $\in$ N such that z$_\text{n}$ lies in S\par
\qfig{p2_math_s3_q02_1}
\qn{3} Let $\vec{\text{a}},\vec{\text{b}},\vec{\text{c}}$ be non-zero non-collinear vectors such that $\vec{\text{c}} = \vec{\text{a}}\times\left(\vec{\text{c}}+\vec{\text{b}}\right)$. The value of $\left|\left(1+\vec{\text{a}}^2\right)\vec{\text{c}}+\vec{\text{b}}\times\vec{\text{a}}+\left(\vec{\text{a}}\times\vec{\text{b}}\right)\times\vec{\text{a}}\right|$ is :
\qn{4} If log$^2$y + $\left(2^{1+\text{x}}+2^{1-\text{x}}\right)$ (log y) + 2$^{2\text{x}+1}$ + 2$^{1-2\text{x}}$ = 0; where x, y $\in$ R, then the value of $|\log(\text{x}+\text{y})|$ is : (Base of log is `e')
\qn{5} Let matrix M = \par\noindent\resizebox{\linewidth}{!}{$\begin{bmatrix}\tan\left(\frac{301\pi}{3}\right) & \sec(2016\pi) & \cot\left(\frac{2015\pi}{2}\right)\\ \cot\left(\tan^{-1}\frac{1}{2}+\tan^{-1}\frac{1}{3}\right) & 2\sin\left(\frac{4\pi}{3}\right) & \sec^{-1}(2016)-\cos^{-1}\left(\frac{1}{2016}\right)\\ \cos\left(\frac{2009\pi}{2}\right) & \cos^{-1}(1) & \sec\left(\frac{301\pi}{3}\right)\end{bmatrix}$}\par , then det.(2M$^\text{T}$ + adj.M) is equal to
\qn{6} Let P be the 7$^\text{th}$ term from the beginning and Q be the 7$^\text{th}$ term from the end in the expansion of $\left(\sqrt[3]{3}+\dfrac{1}{\sqrt[3]{4}}\right)^\text{n}$ where n $\in$ N. If $\dfrac{\text{Q}}{\text{P}}$ = 12, then n is equal to
\end{multicols}
\end{document}
